\section{Introduction}
\label{sec:introduction}

% Scope, organizing principles, functional framework, review approach, and positioning.
% See docs/OUTLINE.md — Scope and organizing principles; Functional framework;
% Review approach; Positioning and novelty.
% Draft 1. Style model: Blumstein (agents/STYLE.md).

Most sounds a rodent hears carry no information about danger, yet failing to respond to the few that do can be fatal. Receivers must therefore decide which sounds matter, how strongly to respond, and when to stop responding to a sound that recurs. Experimental studies show that these decisions are neither fixed nor uniform across receivers. In laboratory rats, selective freezing to conspecific 22-kHz alarm calls depends on prior aversive experience \citep{KimEtAl2010}. Juvenile Richardson's ground squirrels (\textit{Urocitellus richardsonii}) reduce their responses to individual callers whose alarms are repeatedly not followed by a predator \citep{HareAtkins2001},
and rats' approach to playbacks of appetitive 50-kHz calls, strong at first, declines rapidly \citep{WilluhnEtAl2014}. The same acoustic event can thus elicit strong defense, little overt response, or a response that fades, depending on who hears it, what that receiver has experienced, and how often the sound has occurred.

% VERIFY: BlumsteinArmitage1997, SloanEtAl2005, Shriner1998 (not full-text checked): graded information in alarm calls; heterospecific alarm responses. BonautoEtAl2023 and SolajaGustison2026 (not full-text checked): scope of each review.
Two literatures address this variation, but rarely together. Field studies of sciurids, including ground squirrels, marmots, prairie dogs, and chipmunks, have shown that free-living receivers extract graded information from alarm calls \citep{BlumsteinArmitage1997,SloanEtAl2005}, respond to the alarms of other species \citep{Shriner1998}, and weigh the reliability of individual callers \citep{BlumsteinEtAl2004}. Laboratory work on mice and rats has developed playback of ultrasonic vocalizations into a controlled assay of affective state \citep{BonautoEtAl2023} and has begun to measure how such sounds engage neural systems, including striatal dopamine \citep{WilluhnEtAl2014}. The field literature offers natural signal use and ecological costs; the laboratory literature offers stimulus control, developmental manipulation, and access to mechanism. Recent reviews of rodent vocal communication have concentrated on the structure and production of calls \citep{SolajaGustison2026}. How the receiver assigns meaning to a sound, and how that meaning changes with experience and time, has not been synthesized across these systems.

\subsection{A functional framework}
\label{sec:intro-framework}

Four bodies of functional theory make explicit predictions about when receivers should respond and when responsiveness should decline. We use them to frame each question in this review.

First, signal detection theory applied to defense, the ``smoke detector principle,'' holds that when a defense is cheap relative to the cost of missing a real threat, the optimal response threshold produces many false alarms \citep{Nesse2005}. Frequent responses to harmless sounds are therefore expected rather than anomalous. The principle explains why receivers might respond readily to uncertain acoustic cues, but it does not predict which acoustic features they should treat as cues. It frames our treatment of acoustic meaning (\secref{sec:q1}), innate and learned recognition (\secref{sec:q2}), and reliability (\secref{sec:q5}).

% VERIFY: DillHoutman1989 (not full-text checked): flight distance increased with distance to refuge.
Second, threat-sensitive predator avoidance predicts that the intensity of avoidance should scale with the magnitude of the threat \citep{Helfman1989}, and the economics of escape predicts that an animal should flee when the rising cost of staying exceeds the cost of fleeing \citep{YdenbergDill1986}. Gray squirrels (\textit{Sciurus carolinensis}) farther from refuge fled at greater distances from an approaching threat \citep{DillHoutman1989}, although that test used a visual stimulus. In acoustic terms, these models predict that responses should vary with the risk a sound indicates and with the costs of responding, including lost foraging time; graded alarm calls in sciurids provide the relevant natural variation \citep{BlumsteinArmitage1997,SloanEtAl2005}. They frame our analysis of behavioral responses (\secref{sec:q4}).

Third, the predation risk allocation hypothesis predicts that the temporal pattern of risk should shape how much vigilance an animal allocates to each high-risk period: brief, rare pulses of danger warrant intense responses, whereas frequent or prolonged danger forces animals to reduce antipredator effort in order to forage \citep{LimaBednekoff1999}. The hypothesis has not, to our knowledge, been tested with acoustic cues in rodents. We use it as theory that motivates questions about habituation and timing (\secref{sec:q6} and \secref{sec:q7}).

Fourth, receivers can exploit signals not directed at them. Eavesdropping on heterospecific alarm calls may rest on unlearned acoustic similarity or on learning \citep{MagrathEtAl2015}, and the multipredator hypothesis proposes that antipredator responses can persist after the loss of one predator if other predators maintain the same defenses \citep{Blumstein2006}. Eavesdropping is well documented in rodents, but we found no rodent test of the multipredator hypothesis using acoustic cues; it frames the discussion of innate and learned recognition (\secref{sec:q2}).

Together, these frameworks shift the central question from whether a sound elicits a response to when receivers should build, maintain, and revise the meaning they assign to it. Our central claim is that acoustic meaning is not a fixed property of a waveform. It emerges from an interaction among signal structure, receiver predispositions, prior experience, source identity and reliability, environmental state, and the temporal history of exposure.

\subsection{Scope}
\label{sec:intro-scope}

This review concerns the receiver's interpretation of sound, not vocal production. We treat field-studied sciurids and laboratory mice and rats as co-primary systems and note wild-derived or field-experienced laboratory animals where available. Stimuli include conspecific alarm, distress, and social calls; heterospecific alarm calls; predator vocalizations; and synthetic calls, tones, and noise used as acoustic controls. Outcomes include orienting, vigilance, freezing, flight, approach, autonomic and endocrine responses, neural activity, learning, habituation, and recovery. Appetitive conspecific calls, notably rat 50-kHz vocalizations, are included as a contrast: they show that natural, biologically meaningful sounds can engage the same systems with opposite valence, which constrains any claim that a sound is perceived as threatening. Visual, olfactory, and multisensory threat paradigms are discussed only where needed to delimit an acoustic finding. The review is not a survey of sensory fear circuitry or of acoustic deterrent devices.

Throughout, we distinguish acoustic salience, ecological significance, perceived threat, valence, and behavioral deterrence; a startle reflex or a dopamine transient does not by itself indicate sustained avoidance. We also distinguish four levels at which experiments probe receivers: acoustic detection, biological interpretation, action selection, and updating. Because responses at these levels can dissociate, evidence at one level cannot substitute for evidence at another (\secref{sec:synthesis}). For each result, we note whether it was demonstrated directly, inferred from mechanistic data, or remains a hypothesis, and whether receivers were naïve, laboratory-experienced, or field-experienced.

\subsection{Review approach}
\label{sec:intro-approach}

We included studies in which the receiver was a rodent, an acoustic stimulus was experimentally manipulated or naturally varied, and a behavioral, physiological, or neural response was measured. Starting from anchor studies for each question, we traced citations backward and forward and searched for disconfirming evidence for each working claim. Each included experiment is coded by species, setting, receiver experience, stimuli and acoustic controls, endpoints, principal result, and limitation (\tabref{tab:evidence}).
% TODO: search method (databases, strings, dates, counts) deferred until a target journal is chosen; see docs/LITERATURE_MAP.md "Must-haves before drafting".

\subsection{Organization}
\label{sec:intro-organization}

The review is organized as seven questions in three parts. The first part asks what makes a sound biologically meaningful (\secref{sec:q1}) and whether responses to conspecific, heterospecific, and predator sounds are innate or learned (\secref{sec:q2}). The second asks what neural measures reveal about how receivers assess a sound (\secref{sec:q3}) and what determines which behavior a sound evokes (\secref{sec:q4}). The third asks how experience and signal reliability change a sound's meaning (\secref{sec:q5}), whether ecological significance alters habituation and recovery (\secref{sec:q6}), and how the timing of repeated sounds shapes responses (\secref{sec:q7}). We then integrate these questions in a receiver-centered framework (\secref{sec:synthesis}) and derive testable predictions that pair field and laboratory systems (\secref{sec:conclusion}).
